\section{Arithmetic recovery and finite multiplicities}
\label{sec:recovery}

We first work in the weakly mixing setting of Section~\ref{sec:setup}:
\(\Gamma<G\) is an actual ICC lattice, the commuting projective representations
\(\pi,\rho\) on \(K\) are normal and finite over both endpoint factors, and
\(\Ad\rho\) on \(L^2(\pi(\Gamma)')\ominus\mathbb C1\) has no
finite-dimensional subrepresentation.  We then restore the compact sectors and
pass to groups in \(\mathscr K\).

Write \(B_S,B_T\) for the two copies of the full source flag space and
\(m_S,m_T\) for their compact probability measures.  The conclusion of
Section~\ref{sec:gluing} is used in its simultaneous form: the chamber
evaluations are commuting spectral evaluations of nonsingular essentially
invertible maps \(B_S\to B_T\); the evaluations in the two comparisons sharing
slot zero also commute.  The induced spectral measure is equivariant for
source transport and target postcomposition and is fixed by diagonal
\(\rho\).  These simultaneous properties are needed below, not merely the
existence of separate measurable boundary maps.

\subsection{The measurable map space and uninduction}

\begin{lemma}[The space of boundary isomorphisms]
\label{lem:recovery-map-space}
The space \(\mathscr E\) of nonsingular essentially invertible measurable
maps \(B_S\to B_T\), modulo equality almost everywhere, is a standard Borel
space.  Composition with the boundary actions is Borel, and target
postcomposition is continuous for convergence in measure.  The spectral
maps supplied by Section~\ref{sec:gluing} therefore define a PVM on
\(\mathscr E\).
\end{lemma}

\begin{proof}
Equip the compact flag spaces with compatible bounded metrics.  The space
\(L^0(B_S,B_T)\), with distance
\(\int d_T(f(x),g(x))\,dm_S(x)\), is separable and completely metrizable.
There are jointly Borel choices of representatives: approximate a map by
the first member of a countable dense family of simple maps within a
prescribed summable error, take an almost everywhere convergent subsequence
with summable pointwise error bounds, and assign a fixed value where its
pointwise limit does not exist.  Consequently integration of a bounded
Borel function of a map and a source point is Borel in the map.

Choose countable generating algebras \(\mathcal A_S,\mathcal A_T\).
For a map \(f\), write \(m_f=f_*m_S\).  The condition \(m_f\ll m_T\)
is the countable uniform absolute-continuity test
\[
 \forall j\ \exists k\ \forall E\in\mathcal A_T:
 \quad m_T(E)<1/k\ \Longrightarrow\ m_f(E)<1/j.
\]
It is Borel.  The same test with the two measures interchanged recognizes
\(m_T\ll m_f\).  The restriction to a generating algebra suffices, since
that algebra approximates measurable sets in the finite measure
\(m_T+m_f\).

On this equivalence locus, essential invertibility is equivalent to
\[
 \forall E\in\mathcal A_S\ \forall j\ \exists F\in\mathcal A_T:
 \quad m_S(E\mathbin\triangle f^{-1}(F))<1/j.
\]
This is again Borel.  To verify the equivalence, the displayed condition
says that pullbacks generate the completed source sigma algebra.  Recover
each member of a countable separating family of source indicators as a
pullback of a target measurable indicator, using a summable subsequence of
the displayed approximations.  The resulting target indicator sequence
codes a source point almost everywhere for \(m_f\), and hence for \(m_T\).
The inverse of the Borel coding of the source gives a measurable map
\(g:B_T\to B_S\) with \(g f=\id\) almost everywhere.  It follows that
\(f g=\id\) almost everywhere for \(m_f\), hence also for \(m_T\).
Conversely an essential inverse gives the displayed approximations.
Thus \(\mathscr E\) is a Borel subset of \(L^0(B_S,B_T)\).

The jointly Borel representative construction also proves the composition
assertion; nonsingularity makes composition independent of the chosen
representatives.  Target postcomposition is continuous by uniform
continuity on compact flag spaces, also when the acting group element
varies locally.  Source precomposition is Borel by the same integral test.

Here is an explicit countable coding for the spectral pushforward.
Choose a countable real continuous point-separating family $(f_j)$ on
$B_T$ and use the source generating algebra $\mathcal A_S$ above:
\[
 I_{A,j}(\phi)=\int_A f_j(\phi(x))\,dm_S(x),
 \qquad A\in\mathcal A_S.
\]
These coordinates are continuous for convergence in measure; equality of
all of them gives equality of each $f_j\circ\phi$ almost everywhere by
the monotone class theorem, and hence equality of the maps almost
everywhere. The Lusin--Souslin Theorem \cite[Corollary~15.2]{Kechris1995} therefore gives a Borel inverse
to this coding on its image. In the common abelian spectral realization
of Section~\ref{sec:gluing}, the weak integrals
$\int_A P_x(f_j)\,dm_S(x)$ are precisely multiplication by
$I_{A,j}(\phi_\zeta)$ for the jointly measurable random map
$\phi_\zeta$. The chamber-bijection conclusion there puts
$\phi_\zeta\in\mathscr E$ almost surely; faithfulness of the scalar
spectral state therefore puts the joint operator PVM on the coding
image of $\mathscr E$. Its pushforward under the Borel inverse is the
asserted PVM, with covariance and common-slot commutation inherited
from the same integrated operators. To recover a fixed boundary label
$x$, take neighborhoods $U_n$ shrinking to $x$. Source continuity gives
\[
 P_x(f_j)=\mathop{\mathrm{s}\!\!\!-\!\lim}_{n\to\infty}
 \frac{1}{m_S(U_n)}\int_{U_n}P_y(f_j)\,dm_S(y).
\]
Each integral belongs to the coded operator algebra by $L^1$
approximation of source indicators from $\mathcal A_S$. Thus every
fixed-label PVM is recovered, without evaluating an arbitrary
representative of an $L^0$ class.
\end{proof}

The coded PVM acts on \(L^2(Y;K_0\otimes K_1)\), where
\(Y=G/\Gamma\).  The integrated angular operators used in its coding
commute with the base multipliers \(L^\infty(Y)\), so their spectral
projections, and hence the coded PVM, are base-decomposable.
Disintegration over \(Y\) gives PVM kernels \(T_y\) on
\(K_0\otimes K_1\).  Using the same base disintegration for both normal
copies and countable determining families retains their common-slot
commutation outside one null set.  Use the section and cocycle of
Section~\ref{sec:setup}, so
\(c(a,y)=d_0(ay)^{-1}a d_0(y)\).  Its covariance reads
\begin{equation}
 \Ad\pi_0(c(a,y))T_y(f)
   =T_{ay}\bigl(\phi\mapsto f(\phi\circ a)\bigr).
 \label{eq:recovery-induced-covariance}
\end{equation}
The identity holds for each fixed \(a\) outside a null set independent of
\(f\): first use a countable determining algebra and then normality.

\begin{lemma}[Uninduction]
\label{lem:recovery-uninduction}
There is a PVM \(\Theta\) on \(\mathscr E\), acting on
\(K_0\otimes K_1\), such that, for every \(g,h\in\Gamma\),
\begin{align}
 \Ad\pi_0(g)\Theta(f)
   &=\Theta\bigl(\phi\mapsto f(\phi\circ g)\bigr),\label{eq:recovery-trunk}\\
 \Ad\pi_1(h)\Theta(f)
   &=\Theta\bigl(\phi\mapsto f(h^{-1}\circ\phi)\bigr).
   \label{eq:recovery-branch}
\end{align}
It is fixed by \(\Ad(\rho_0(\lambda)\otimes\rho_1(\lambda))\)
for every \(\lambda\in\Lambda\).
The two copies of \(\Theta\) sharing slot zero commute.
\end{lemma}

\begin{proof}
For \(g=d_0(y)h\) define the lifted kernel
\[
 F_g(f)=\Ad\pi_0(h)^{-1}
     T_y\bigl(\phi\mapsto f(\phi\circ g)\bigr).
\]
The tests depending on \(g\) are legitimate Borel kernel tests by
Lemma~\ref{lem:recovery-map-space}; approximate them by simple Borel tests
when applying the kernel identities.  Since
\(ag=d_0(ay)c(a,y)h\), Equation~\eqref{eq:recovery-induced-covariance}
gives \(F_{ag}(f)=F_g(f)\).  All scalar matrix coefficients of a countable determining family are
measurable and invariant under each left translation. Haar invariance
and Fubini make them essentially constant. Choose one $g_0$ in the
common conull set on which these coefficients have their constant
values and the lifted kernel is a PVM, and define $\Theta=F_{g_0}$.
For almost every $g$, equality on the generating algebra and uniqueness
of the finite scalar matrix-coefficient measures give $F_g=\Theta$
on all Borel sets. Choose the same good $g_0$ for both normal copies
and their countable determining commutators. This retains their
common-slot commutation. Thus $\Theta$ is a PVM by construction;
its countable additivity is inherited from that fibre.

Replacing \(g\) by \(g\gamma\) in the lift gives
\[
 F_{g\gamma}(f)=\Ad\pi_0(\gamma)^{-1}
   F_g\bigl(\phi\mapsto f(\phi\circ\gamma)\bigr),
\]
which proves Equation~\eqref{eq:recovery-trunk}.
Target covariance gives Equation~\eqref{eq:recovery-branch}.
The diagonal \(\rho\)-invariance and common-slot commutation hold for the
induced kernels.  Lifting the two comparisons simultaneously and taking
the same countable determining families preserves those identities.
\end{proof}

\begin{lemma}[A target transversal]
\label{lem:recovery-transversal}
Target postcomposition of \(\Gamma\) on \(\mathscr E\) is free and admits
a Borel transversal.
\end{lemma}

\begin{proof}
If \(h\circ\phi=\phi\) almost everywhere, equivalence of
\(\phi_*m_S\) and \(m_T\) implies that \(h\) fixes almost every target
flag.  Continuity and full support make it fix every flag, and faithfulness
of the effective flag action gives \(h=e\).

We prove the local slice property.  Suppose
\(\phi_n\to\phi\) and \(h_n\circ\phi_n\to\psi\) in measure in
\(\mathscr E\).  An escaping sequence \(h_n\) has, after a subsequence,
convergent compact Cartan frames and at least one diverging simple gap.
In the corresponding proximal vertex module its normalized action
contracts compact sets outside a limiting proper algebraic repeller to a
single vertex.  These are the local-field contraction properties of
Section~\ref{sec:setup}, applied to one factor if necessary.
Pass to almost everywhere convergence of both map sequences.  Since
\(\phi_*m_S\sim m_T\), its values avoid the repeller almost everywhere.
At each such point, \(\phi_n\) is eventually in a compact neighborhood
away from the repeller, and the chosen vertex coordinate of
\(h_n\phi_n\) tends to the attracting vertex.  The same coordinate of
\(\psi\) is therefore constant almost everywhere, contradicting
\(\psi_*m_S\sim m_T\).  Thus the translating sequence cannot escape.

For each \(\phi\) there is consequently a neighborhood \(U\) with
\(hU\cap U=\varnothing\) for every \(h\ne e\).  Otherwise select pairs
in progressively smaller neighborhoods with nonidentity translators.
Escaping translators have just been excluded, and a constant subsequence
would contradict freeness by continuity.  Choose a countable cover by
such neighborhoods, using separability.  Each member meets an orbit at
most once.  Retain from the first member all its points, and from each
successive member remove the countable saturation of the preceding
members.  The resulting Borel set meets each orbit exactly once.
\end{proof}

\subsection{From a simultaneous boundary law to a fixed partition}

We now use the PVM $\Theta$ obtained by uninduction.  Its two comparison
copies have a common first module slot.  Their commutation will let us
regard the law in the second slot as a measurable family of ordinary
PVMs.  Each member of that family gives a finite-rank partition of $K$.
The point of the argument is to show that these partitions are the same,
even though their boundary-map labels can vary.

We first record the consequence of normality used to compare partitions.

\begin{lemma}[Normal coefficient decay]
\label{lem:recovery-coefficients}
If a projective group representation extends normally to its twisted
group factor, its matrix coefficients vanish at infinity.  In particular,
for finite-rank projections $p,q$,
\[
 \Tr\bigl(p\pi(g)q\pi(g)^*\bigr)\longrightarrow0
 \qquad(g\longrightarrow\infty).
\]
\end{lemma}

\begin{proof}
A normal coefficient is given by an $L^1$-density on the finite factor.
Approximate that density in $L^1$ by a finite Fourier polynomial.  The
polynomial coefficient at $u_g$ vanishes outside a finite set, and the
approximation error is uniformly bounded by the $L^1$-error.  Expanding
$p,q$ in finite orthonormal bases gives the displayed conclusion.
\end{proof}

\begin{proposition}[A deterministic atom partition]
\label{prop:recovery-partition}
Retain the weakly mixing correspondence $K$ from the start of this section.
Let $\Theta$ be a PVM of total mass $I$ on the boundary-isomorphism space
$\mathscr E$, acting on $K_0\otimes K_1$, with the precomposition and postcomposition
covariances of Equations~\eqref{eq:recovery-trunk} and
\eqref{eq:recovery-branch}.  Suppose that diagonal $\rho$ fixes $\Theta$
and that its two copies on $K_0\otimes K_1\otimes K_2$ commute.
Then $K$ has an orthogonal partition into nonzero finite-dimensional
subspaces permuted by $\pi(\Gamma)$ and $\rho(\Lambda)$.  The
$\Gamma$-action on the partition is free and has finitely many orbits.
\end{proposition}

\begin{proof}
\emph{Disintegration and finite multiplicity.}
Take all normal slot-one slices of $\Theta(f)$ and let $\mathcal A_0$
be the von Neumann algebra they generate on $K_0$.  Apply two normal
functionals, one in each branch slot, to the common-slot commutation
identity.  The resulting slices commute, also after taking adjoints;
thus $\mathcal A_0$ is abelian.  The covariance identities normalize it
by $\pi_0$ and $\rho_0$.  The slice criterion gives
\[
 \Theta(f)\in\mathcal A_0\bar\otimes B(K_1).
\]

Use a faithful normal state to represent $\mathcal A_0$ over a standard
probability space $Z$.  The displayed inclusion disintegrates $\Theta$
into PVMs $\Theta_z$ on the fixed Hilbert space $K_1$.  This can be done
first for one bounded real Borel coordinate generating $\mathscr E$ and
then for its spectral calculus.  Countably many generators and the
countable group $\Gamma$ make target covariance hold simultaneously
outside one null set.

Let $T\subset\mathscr E$ be the Borel target transversal from
Lemma~\ref{lem:recovery-transversal}, and set $q_z=\Theta_z(T)$.  The
projections $\pi_1(g)q_z\pi_1(g)^*$, $g\in\Gamma$, are orthogonal
and sum to one.  Hence
\[
 \ell^2(\Gamma)\otimes q_zK_1\longrightarrow K_1,
 \qquad \delta_g\otimes\eta\longmapsto\pi_1(g)\eta
\]
is a unitary intertwining the first-factor twisted regular action with
$\pi_1$.  In particular,
\[
 \dim_{L_\mu(\Gamma)}K=\dim_{\mathbb C}q_zK_1<\infty.
\]
The PVM on $q_zK_1$ therefore has finitely many nonzero atoms.  Their
$\Gamma$-translates give all of $\Theta_z$, so its support consists of
finitely many free target orbits of finite-rank atoms.

\emph{A measurable encoding of the projection partition.}
Choose a strictly positive trace-class operator $D$ on $K_1$.  The
measures
\[
 E\longmapsto\Tr\bigl(D\Theta_z(E\cap T)\bigr)
\]
form a measurable kernel.  Their nonzero atoms are exactly the nonzero
atoms of the transversal PVM.  After a Borel injection of $\mathscr E$
into $[0,1]$, enumerate these atoms in increasing order: the first atom
is the infimum of the rational cuts with positive mass; remove that
mass and repeat.  This is measurable and stops after at most
$\dim q_zK_1$ steps.  Kernel integration over the resulting Borel graphs
makes the singleton projection weights measurable.  Their matrix
coefficients and finite ranks make them measurable also as elements of
the separable Hilbert--Schmidt space.

Let $\mathcal P$ be the space of nonzero finite-rank orthogonal
projections on $K_1$.  On its $\Ad\pi_1(\Gamma)$-orbit space put
\[
 d_\Gamma([p],[q])
   =\inf_{g\in\Gamma}
       \|p-\pi_1(g)q\pi_1(g)^*\|_{\mathrm{HS}}.
\]
The triangle inequality follows from the isometric action.  To see that
the distance separates orbits, note that
\[
 \|p-\pi_1(g)q\pi_1(g)^*\|_{\mathrm{HS}}^2
 =\rank p+\rank q-2\Tr\bigl(p\pi_1(g)q\pi_1(g)^*\bigr).
\]
Distance zero requires equal ranks.  Lemma~\ref{lem:recovery-coefficients}
excludes an escaping minimizing sequence, while a nonescaping
subsequence gives equality after one translation.  The orbit metric is
therefore separating and separable.  Its nonempty finite subsets carry
the separating, separable Hausdorff metric.  These distances are
measurable in the preceding enumeration, since every infimum is over
the countable group $\Gamma$.

For each $z$, let $P(z)$ be the finite set of orbit classes of the
transversal atom projections.  This encoding forgets the boundary-map
labels and the choice of a representative projection in each target
orbit.  It retains the full projection partition: take the union of the
$\Gamma$-orbits of the projections represented by $P(z)$.
Source precomposition relabels the boundary maps and commutes with
target postcomposition. Returning a moved map to the chosen transversal
$T$ only conjugates its projection weight by $\pi_1(\Gamma)$.
Thus source precomposition leaves $P(z)$ unchanged, and every bounded
measurable function of $P(z)$ commutes with $\pi_0$.  Let $\mathcal C_0$ be the resulting
abelian algebra.  Then
\[
 \mathcal C_0\subset\pi_0(\Gamma)'.
\]

\emph{Weak mixing removes the remaining variation.}
The normalized trace on the finite factor $\pi_0(\Gamma)'$ gives the
distribution of $P$ an invariant probability measure.  It has the same
null projections on $\mathcal C_0$ as the faithful normal state used to
disintegrate $\mathcal A_0$.  A deterministic conclusion for this trace
therefore also holds in the original spectral realization.
Diagonal $\rho$-invariance of $\Theta$ gives, for bounded functions $F$,
\[
 \Ad\rho_0(h)F(P)
    =F\bigl(\Ad\rho_1(h)^{-1}P\bigr).
\]
Because $\rho_1$ commutes with $\pi_1$, it acts isometrically on the
projection-orbit space and on its finite-subset space.

Weak mixing of $\Ad\rho$ on
$L^2(\pi_0(\Gamma)')\ominus\mathbb C1$ makes the diagonal action on
$\mathcal C_0\bar\otimes\mathcal C_0$ ergodic.  Indeed, a nonconstant
invariant tensor with a mean-zero factor gives a nonzero Hilbert--Schmidt
intertwiner.  The associated positive compact operator would have a
nonzero finite-dimensional invariant spectral subspace, contradicting
weak mixing.  The bounded truncation of the distance between two
independent values of $P$ is thus constant almost everywhere.  That
constant is zero: separability gives a countable cover by balls of
arbitrarily small radius, and a ball of positive probability gives a
positive probability of arbitrarily small pair-distance.  Fubini now
makes $P$ itself constant almost everywhere.

The full projection partition recovered from this constant value is
fixed on $K_1=K$.  Covariance makes it invariant as a partition under
both endpoint groups.  The free $\Gamma$-action and its finite orbit
count are those already obtained from the transversal decomposition.
\end{proof}

\subsection{From a finite-rank partition to arithmetic models}

The remaining group-theoretic step no longer uses the boundary law or
weak mixing.  We state its actual partition hypothesis so that the same
result can later be applied after restoring the compact sectors.

\begin{proposition}[Arithmetic models from a partition]
\label{prop:recovery-graph-model}
Let $\Gamma,\Lambda$ be countably infinite ICC groups, with normalized
scalar cocycles $\mu,\omega$, and write $M=L_\mu(\Gamma)$ and
$N=L_\omega(\Lambda)$.  Let $K$ be a nonzero normal bifinite
$M$--$N$ correspondence.  Suppose that
$K$ has an orthogonal partition into nonzero finite-dimensional subspaces
permuted by $\pi(\Gamma)$ and $\rho(\Lambda)$, and that $\Gamma$ acts
freely on its labels with finitely many orbits.  Then $K$ is a finite
orthogonal sum of the models $E(A,B,\delta,\sigma)$ of
Theorem~\ref{thm:arithmetic-exhaustion}.
\end{proposition}

\begin{proof}
Let $X$ be the label set and fix $x\in X$.  Its stabilizer in $\Lambda$
is finite: an infinite stabilizer would preserve the corresponding
finite-rank projection along an escaping sequence, contradicting
Lemma~\ref{lem:recovery-coefficients}.  Since there are finitely many
$\Gamma$-orbits, a finite-index subgroup of $\Lambda$ preserves
$\Gamma x$.  The two permutation actions commute, so moving $x$ inside
$\Gamma x$ does not change its $\Lambda$-stabilizer.  This finite-index
subgroup therefore normalizes that finite stabilizer.  Each element of
the finite stabilizer has a finite conjugacy class in $\Lambda$, and ICC makes the
stabilizer trivial.

Thus $\Lambda$ also acts freely.  On each of its orbits, a basis of one
atom gives that many copies of the twisted regular right module.  Finite
right module dimension makes the number of these orbits finite, so there
are finitely many joint orbits as well.

Fix a joint orbit and its base atom $V_0$.  Let $A\leq\Gamma$ consist
of the elements taking its label into its $\Lambda$-orbit, and define
$B\leq\Lambda$ symmetrically.  Their indices are finite by the two
individual orbit counts.  Separate freeness gives, for each $a\in A$,
a unique $\delta(a)\in B$ such that $\pi(a)\rho(\delta(a))$ preserves
$V_0$.  The joint stabilizer is a subgroup of
$\Gamma\times\Lambda$ whose coordinate projections are injective.
It is consequently the graph of an actual isomorphism $\delta:A\to B$.
Its action on $V_0$ is
\[
 \sigma(a)=\left.\pi(a)\rho(\delta(a))\right|_{V_0},\qquad
 \sigma(a)\sigma(b)
   =\frac{\mu(a,b)}{\omega(\delta(a),\delta(b))}\sigma(ab).
\]

To identify the joint-orbit span with its arithmetic model, first define
\[
 F_0:V_0\otimes L^2(N)\longrightarrow K,
 \qquad F_0(\xi\otimes v_h)=\xi v_h.
\]
Orthogonality of the distinct $\Lambda$-translates makes $F_0$ a
right-module isometry onto their closed span.  The graph stabilizer gives
\[
 u_a\xi=\sigma(a)\xi\,v_{\delta(a)},
\]
so $F_0$ intertwines the stipulated left subgroup action.  This subgroup
action is normal by projective regular absorption
(Lemma~\ref{lem:setup-twisted-absorption}) on each $B$-orbit in $\Lambda$.

Choose representatives $r\in\Gamma/A$.  As a right subgroup module,
\[
 L^2(M)=\bigoplus_r u_rL^2(L_{\mu|A}(A)).
\]
The relative tensor product defining $E(A,B,\delta,\sigma)$ is therefore
the orthogonal sum of the spaces
$u_r\otimes(V_0\otimes L^2(N))$.  On them set
\[
 u_r\otimes\eta\longmapsto\pi(u_r)F_0(\eta).
\]
Distinct cosets give distinct $\Lambda$-orbits of atoms, making this map
isometric with range the full chosen joint-orbit span.  Balancing and
the subgroup identity give both endpoint intertwining identities on
finite Fourier tensors, and hence on the completion.  This is the
required unitary with $E(A,B,\delta,\sigma)$.  Summing over the finitely
many joint orbits proves the proposition.
\end{proof}

\begin{proposition}[Arithmetic exhaustion]
\label{prop:recovery-exhaustion}
For every $\Gamma\in\mathscr K$, every countably infinite ICC
$\Lambda$, and arbitrary scalar cocycles $\mu,\omega$, every nonzero
separable bifinite $L_\mu(\Gamma)$--$L_\omega(\Lambda)$ correspondence
is a closed summand of a finite direct sum of the arithmetic models in
Theorem~\ref{thm:arithmetic-exhaustion}.
\end{proposition}

\begin{proof}
Restrict the source to a finite-index subgroup actually isomorphic to a
finite-index subgroup of a lattice, with its exactly pulled-back
cocycle.  The latter subgroup is itself an ICC lattice.  By
Lemma~\ref{lem:setup-compact-reduction}, a finite-index restriction of
$\Lambda$ splits the correspondence into finitely many blocks
\[
 K_z=V_c\otimes K',\qquad
 \pi=1\otimes\pi',\qquad
 \rho_h=v_h\otimes\rho'_h.
\]
If $v$ has multiplier $\nu$, then $\rho'$ has right-adjoint multiplier
$\overline\omega/\nu$, and its right factor has cocycle $\omega\nu$.
The reduced correspondence is normal, bifinite, and satisfies the weak
mixing hypothesis used in the geometric construction and in
Proposition~\ref{prop:recovery-partition}.

Apply that proposition to $K'$ and tensor its fixed atom partition with
$V_c$.  Both original actions permute the resulting finite-dimensional
atoms.  Their label actions are unchanged, so the $\Gamma$-action is
still free with finitely many orbits.  The hypotheses of
Proposition~\ref{prop:recovery-graph-model} now hold for the original
restricted correspondence, without a weak mixing assumption on its
restored actions.  It gives the arithmetic models with the original
endpoint phases.  Explicitly, on a recovered graph the reduced fibre
multiplier is
$\mu/(\delta^*\omega\,\delta^*\nu)$; tensoring with
$a\mapsto v_{\delta(a)}$ multiplies it by $\delta^*\nu$ and restores
$\mu/\delta^*\omega$.  No finite-image assumption on $v$ is used.

Finally apply Lemma~\ref{lem:setup-finite-index-induction} on both sides.
The original correspondence embeds as a closed summand of the two-sided
induction of its restriction.  The normalized embeddings are the coset
sums with terms
\[
 u_r\otimes\pi(u_r)^*\xi,
 \qquad
 (\xi v_s^*)\otimes v_s,
\]
respectively.  Associativity of relative tensor products and the standard
trace-module identifications show that inducing each restricted model
gives
\[
 L^2(M)\otimes_{L_{\mu|A}(A)}
       \bigl(V_\sigma\otimes L^2(N)\bigr),
\]
with the same graph subgroup, viewed now in the ambient groups.  All
indices and sums are finite, proving exhaustion.
\end{proof}

\begin{proof}[Completion of Theorem~\ref{thm:arithmetic-exhaustion}]
Proposition~\ref{prop:recovery-exhaustion} gives the forward assertion
for every nonzero correspondence; the zero correspondence is the zero
summand.  Corollary~\ref{cor:models-dimensions} gives normality and both
dimensions of each elementary model.  Normality and finite module
dimension pass to finite direct sums and closed bimodule summands, as
proved at the end of Section~\ref{sec:models}.  This supplies the converse
and completes the theorem.
\end{proof}
